\section{Methods}\label{sec:methods}

\subsection{Data Source, Record Audit and Common Support}\label{sec:data}

We conducted a cross-sectional analysis of outbreak-level records in the public CDC NORS dataset~\cite{cdc2024norsdata}, downloaded on 25 November 2026 (portal rows last updated 20 December 2024), adhering to RECORD reporting standards~\cite{benchimol2015record} (full checklist in Supplementary Section S4 / Table S12). Across NORS reporting revisions spanning 2009--2023~\cite{cdc2014norsguidance,wikswo2022enteric}, the core variables---estimated total primary ill persons (\texttt{Illnesses}, $C_i$), hospitalizations ($H_i$), and persons with known hospitalization status ($K_i$)---retain consistent definitions~\cite{cdc2014norsguidance,cdc2024norsdata}. As documented in CDC NORS guidance~\cite{cdc2014norsguidance} (pp.~15--16), $C_i$ combines total ill persons ``including lab-confirmed and probable cases'' based on case definitions or estimates, so $M_i = C_i - K_i$ mixes individually ascertained cases lacking hospitalization records with estimated cases (analogous to ``missing types''~\cite{jiang2018using}); because the public NORS extract does not separate confirmed, probable and estimated counts, missing-type bound-tightening would require restricted data. Reporting by health departments is voluntary~\cite{hall2013acute}, and etiologies are classified as confirmed or suspected~\cite{cdc2026surveillance2019}. Eligible primary outbreaks had food as primary mode, began in 2009--2023, and had confirmed etiology in one of six study groups: any \textit{Salmonella} serotype, STEC, any \textit{Campylobacter} species, \textit{C. perfringens}, norovirus GI or norovirus GII (single or same-group multi-entry~\cite{cdc2014norsguidance}; single-entry~\cite{weller2026single}, confirmed O157-restricted STEC~\cite{tack2021shiga} [279/388 outbreaks], and excluding 6 \textit{Salmonella} serovar Typhi outbreaks [87 cases, 42 hospitalized] in sensitivity analyses). Excluded groups comprised 951 suspected study-group outbreaks (included in variant R8), 304 confirmed norovirus outbreaks lacking genogroup (only 2 recoverable from recorded genotype tokens; included in R7b), and 300 mixed/non-study outbreaks (104 with $\ge 1$ confirmed study group; Supplementary Table S2).

Complete-field consistent records satisfied $0 \le H_i \le K_i \le C_i$ with both $H_i$ and $K_i$ recorded (3,976 outbreaks), whereas \textit{fully recorded outbreaks} denote zero unrecorded cases ($M_i = 0$; 2,414 outbreaks in common support). Incomplete records were retained under logical bounds (Supplementary Table S1): blank $K_i$ set $K_{\text{eff},i} = H_i$ (leaving $C_i - H_i$ unrecorded); blank $H_i$ with $K_i = 0$ implied $H_i = 0$; and blank $H_i$ with $K_i > 0$ (23 outbreaks), or both fields blank, set $K_{\text{rec},i} = 0$ and treated all $C_i$ cases as unrecorded (in the 23 blank-$H$, $K_i > 0$ outbreaks, this discards $H_i \le K_i$ harmlessly for worst-case validity while preventing point-rate distortion). Contradictions ($H_i > C_i$ or $K_i > C_i$) or $C_i \le 0$ were excluded.

Seven harmonized settings---restaurant, catering/banquet, private home, institution, other, multiple and unknown (Supplementary Section S1)---were crossed with three periods (2009--2013, 2014--2018, 2019--2023). The primary analysis restricted attention to the data-selected common support of cells with $\ge 3$ outbreaks per group and directly standardized all six groups~\cite{keiding2014standardization} to pooled-case weights $w_s = C_{\cdot s} / \sum_{s'} C_{\cdot s'}$. Under pooled-case weighting, the standard population draws 72.8\% (52,507 of 72,132 cases) from \textit{Salmonella} and norovirus GII, so $\theta_j$ reflects hospitalization under that empirical case mix. All inference is conditional on this selected support; we also evaluated support thresholds $\ge 2$ (R4a) and $\ge 5$ (R4b), period-only standardization (R2), equal within-cell outbreak weighting (R10, normalizing both outbreak $H_i/C_i$ and known-status $H_{\text{rec},i}/K_{\text{rec},i}$ counts per outbreak), alternative external weights (R6a--c), healthcare/LTCF exclusion (R9), and single-state-only restriction (R11).

\subsection{Target Estimand, Worst-Case Bounds and Breakdown Frontiers}\label{sec:bounds}

For group $j$ and cell $s$, let $C_{js}$, $H_{js}$, $M_{js}$, $H_{\text{rec},js}$, and $K_{\text{rec},js}$ denote cell sums of $C_i$, $H_i$, $M_i$, and recorded-status counts ($H_{\text{rec},i} = H_i, K_{\text{rec},i} = K_{\text{eff},i}$ for consistent records; $0, 0$ when $H_i$ is blank and $K_i > 0$ or both are blank). The target estimand is the case-weighted directly standardized hospitalization proportion among \textit{reported foodborne outbreak cases}:
\begin{equation}
  \theta_j = \sum_s w_s \, \frac{H_{js} + H^{u}_{js}}{C_{js}},
  \label{eq:theta}
\end{equation}
where $H^{u}_{js} \in [0, M_{js}]$ is unobserved hospitalizations among unrecorded cases. All bounds take the estimated case counts $C_i$ and reported counts $(H_i, K_i)$ as exact; differential case ascertainment is not bounded here (\S\ref{sec:limitations}). Without assumptions on unrecorded hospitalization outcomes~\cite{manski1989anatomy}, $\theta_j$ lies in $[L_j, U_j]$:
\begin{equation}
  L_j = \sum_s w_s \frac{H_{js}}{C_{js}}, \qquad
  U_j = \sum_s w_s \frac{H_{js} + M_{js}}{C_{js}}.
  \label{eq:bounds}
\end{equation}
Pair $(h, l)$ separates in the sample without outcome assumptions if $L_h > U_l$~\cite{horowitz2000nonparametric}. Let $p_{js} = H_{\text{rec},js}/K_{\text{rec},js}$ be the within-cell recorded rate ($\delta = 0$ assumes within-cell missingness at random, MAR). Motivated by bounded-variation partial identification~\cite{manski2018how}, we evaluated:
\begin{enumerate}
  \item \textit{Additive bounded-variation bands and breakdown points:} allowing each cell's unrecorded rate in $[p_{js} - \delta_{\text{down}},\, p_{js} + \delta_{\text{up}}] \cap [0, 1]$~\cite{ratitch2013missing,liublinska2014sensitivity}:
  \begin{equation}
  \begin{aligned}
    L_j(\delta_{\text{down}}) &= \sum_s w_s \frac{H_{js} + M_{js}\max(0,\, p_{js} - \delta_{\text{down}})}{C_{js}}, \\
    U_j(\delta_{\text{up}}) &= \sum_s w_s \frac{H_{js} + M_{js}\min(1,\, p_{js} + \delta_{\text{up}})}{C_{js}},
  \end{aligned}
    \label{eq:delta}
  \end{equation}
  where symmetric $\delta_{\text{down}} = \delta_{\text{up}} = \delta$ recovers worst-case $[L_j, U_j]$ at $\delta = 100\text{ pp}$. For the 52 blank-$K_i$ outbreaks (2.3\% of eligible cases; true $K_i \in [H_i, C_i]$), setting $K_{\text{eff},i} = H_i$ is exact for worst-case bounds and conservative for the $\delta$-upper bound; a strictly conservative $\delta$-lower bound holding $H_i$ fixed is verified in sensitivity checks. Sample breakdown points $\delta^{*}_{hl} = \inf\{\delta \ge 0 : L_h(\delta) \le U_l(\delta)\}$~\cite{kline2013sensitivity,masten2020inference} were solved analytically (Supplementary Section S2), alongside 2D breakdown frontiers $(\delta_{h,\text{down}}, \delta_{l,\text{up}})$ and asymmetric $-10/+5, -15/+5, -20/+5\text{ pp}$ plus diagnostic-informed group-specific asymmetric bands (\textit{Salmonella} $-10/+5\text{ pp}$, \textit{Campylobacter} $-7/+5\text{ pp}$, others $\pm 5\text{ pp}$). We use the additive scale as our primary target because public-health policy and case-burden comparisons depend on absolute hospitalization risk differences rather than relative odds ratios across groups spanning 0.5\% to 28.5\%; we anchor $\delta$ externally against US per-case benchmarks~\cite{scallan2011foodborne,scallanwalter2025foodborne} and low-missingness states, and also check outbreak-level MAR (R13) and common-mechanism scaling $q_{js} = \min(1, \alpha p_{js})$ (Supplementary Fig.~S4).
  \item \textit{Odds-ratio (OR) pattern-mixture scale:} bounding unrecorded odds in $[\text{odds}(p_{js})/\Gamma,\, \Gamma\cdot\text{odds}(p_{js})]$ ($\Gamma \ge 1$) with sample critical breakdown $\Gamma^{*}_{hl}$, alongside a $0.1\%$ continuity floor ($p_{js} = 0.001$) on 9 zero-rate cells and formal 30-hypothesis inference ($\Gamma = 1.5, 2, 3, 5$ and $\Gamma^{*}_{0.95}$).
  \item \textit{Denominator contamination and selection sensitivity:} drawing on contaminated-data bounds~\cite{horowitz1995identification,molinari2008partial}, we added variants treating 25\% of unrecorded cases as extrapolated non-cases ($C'_i = C_i - 0.25 M_i$, R14a) and a selection-sensitivity variant restricting to outbreaks with $\le 50\%$ unrecorded status ($M_i/C_i \le 0.50$, R14b).
  \item \textit{Composition-adjusted partial vs.\ full diagnostics:} comparing known-status rates between partially ($M_i > 0, K_{\text{rec},i} > 0$) and fully ($M_i = 0$) recorded outbreaks crudely, across shared cells containing both partially and fully recorded outbreaks (within-cell difference and common-cell standardized observed-minus-full difference, $\theta^{\text{cc}}_j(0) - r^{\text{cc}}_{\text{full},j}$, evaluated over identical shared cells and weights), stratified by single-state vs.\ multistate status, and descriptively in 11 high-completeness states satisfying both $<5\%$ unrecorded cases and $\ge 200$ common-support cases (Wilson 95\% CIs).
\end{enumerate}

\subsection{Superpopulation Target, 30-Hypothesis Inference and Double-Bootstrap Evaluation}\label{sec:inference}

Although NORS approaches a national census of reported outbreaks in a fixed historical window, statistical inference treats the observed outbreaks in each cell as an independent stratified draw from a superpopulation of foodborne outbreak-generation and reporting conditional on the selected common support~\cite{keiding2014standardization}. This working assumption ignores cross-outbreak correlation (which we examine directly via exposure-state clustering, CV3 jackknife, and Bayesian cluster resampling below); hierarchical shrinkage~\cite{normand1997statistical} is a complementary approach for sparse cells.

For all $30$ one-sided hypotheses $H_{0,(h,l)}: \theta_h \le \theta_l$ (both directions of all 15 pairs, avoiding data-selected testing~\cite{andrews2024inference} and framing inference on many one-sided moment inequalities~\cite{chernozhukov2019inference}), each lower difference $\Delta^{L}_{hl} = L_h - U_l$ is smooth away from kinks at $p_{js} = \delta$ or $1 - \delta$ (0 of 72 cells lay within $10^{-6}$ of kinks at $\delta = 5\text{ pp}$; minimum distance across all finite 95\% roots was $0.0013\text{ pp}$ under CV3, $0.0024\text{ pp}$ under HC1, and $0.0056\text{ pp}$ under bootstrap-$t$; estimators at kinks can exhibit non-Gaussian directional limits~\cite{fang2019inference,masten2020inference}). Outbreak-HC1 sandwich covariances $\widehat{V}_{\text{HC1}}$ (Eq.~\ref{eq:sandwich}) matched closed-form influence functions $\psi_i$ to $\le 1.64 \times 10^{-12}$ (worst case), $2.92 \times 10^{-13}$ ($\delta = 5\text{ pp}$), and $4.00 \times 10^{-13}$ ($-15/+5\text{ pp}$), while OR ($\le 6.40 \times 10^{-13}$) and group-specific asymmetric ($\le 3.16 \times 10^{-13}$) influence functions matched central-difference covariances and vectorized bootstrap-$t$ expansions agreed to $\le 5.6 \times 10^{-17}$ on $\Delta^L$ and $\le 6.9 \times 10^{-18}$ on standard errors:
\begin{equation}
  \widehat{V}_{\text{HC1}} = \sum_{j=1}^{G}\sum_{s=1}^{K} \widehat{J}_{js}\, \widehat{\Omega}_{js}^{\text{HC1}}\, \widehat{J}_{js}^{\top}, \qquad
  \widehat{\Omega}_{js}^{\text{HC1}} = \frac{1}{n_{js}(n_{js} - 1)} \sum_{i \in (j, s)} (Z_i - \bar{Z}_{js})(Z_i - \bar{Z}_{js})^{\top}.
  \label{eq:sandwich}
\end{equation}
Because NORS's exposure-state field comprises 52 single jurisdictions plus a CDC-coordinated Multistate category (219 outbreaks), we formed 53-cluster exposure-state CV1 covariances $\widehat{V}_{\text{state}}$, computed each cluster's share of total score variance $s_g = \|\Phi_g\|_2^2 / \sum \|\Phi_{g'}\|_2^2$ and plug-in effective cluster count $G_{\text{eff}} = 1/\sum s_g^2$~\cite{carter2017asymptotic,mackinnon2023cluster} (noting that plug-in $G_{\text{eff}}$ from stacked influence scores approximates the scalar variance-ratio effective cluster count of~\cite{carter2017asymptotic,mackinnon2023cluster}), and added three cluster-robust checks less vulnerable to a dominant cluster: (i) leave-one-cluster-out \textbf{CV3 jackknife} variances~\cite{mackinnon2023cluster} (all 53/53 and 52/52 leave-cluster-out tables valid), (ii) a \textbf{Bayesian 53-cluster bootstrap-$t$} drawing i.i.d.\ $\text{Exp}(1)$ cluster weights ($B = 2,000$), and (iii) splitting Multistate into 219 singleton clusters (271 clusters) alongside full inference within the 52 single states of variant R11.

Recognizing that severe outbreak skewness and small cells ($N_{\text{eff}} \approx 2$) distort asymptotic Gaussian critical values and standard cluster approximations~\cite{cameron2008bootstrap,mackinnon2017wild,canay2021wild,mackinnon2023cluster}, we evaluated six 30-hypothesis stepdown Max-$T$ procedures~\cite{romano2005exact} without claiming formal pre-specification: (1) outbreak-HC1 Gaussian (200,000 multiplier draws), (2) 53-cluster exposure-state CV1 Gaussian, (3) 53-cluster CV3 jackknife Gaussian, (4) 9,999-draw unrestricted Rademacher wild cluster bootstrap (WCB), (5) Bayesian 53-cluster bootstrap-$t$ ($B = 2,000$), and (6) outbreak-stratified replicate-studentized bootstrap-$t$ ($T^{(b)}_m = (\Delta^{L,(b)}_m - \widehat{\Delta}^L_m)/\widehat{\sigma}_{\text{HC1},m}^{(b)}$, $B = 2,000$, zero degenerate SE entries). Pointwise Imbens--Manski~\cite{imbens2004confidence} 95\% intervals are also reported with explicit finite-sample under-coverage caveats~\cite{stoye2009more}.

Simultaneous 95\% rank confidence sets $[1 + m^{-}_j,\, G - m^{+}_j]$ count the number of competitors rejected significantly above ($m^{-}_j$) and significantly below ($m^{+}_j$) group $j$ at family-wise $\alpha = 0.05$; these rank intervals are valid and conservative under partial identification because $\theta_h - \theta_l \ge \Delta^L_{hl}$ across the identified set~\cite{klein2020joint,mogstad2024inference} (partitioning methods~\cite{almohamad2021simultaneous,almohamad2022simultaneous} can shorten intervals for small $G = 6$ under point identification). Following~\cite{kline2013sensitivity,masten2020inference}, scalar statistical breakdown points $\delta^{*}_{0.95}$ and $\Gamma^{*}_{0.95}$ are inverse one-sided lower confidence bounds tested pointwise in $\delta$ or $\Gamma$ using fixed-level critical values ($c(\delta)$ from a 21-point grid for Gaussian procedures and 14-point grid for $c_{\text{bt}}(\delta)$, with all finite bootstrap-$t$ roots refined by exact uninterpolated $c_{\text{bt}}(\delta)$ evaluation to residual $\le 3.1 \times 10^{-5}$ and unique crossings verified on a 201-point scan; simultaneous confidence bands over the 2D frontiers in Fig.~4 would require a uniform band~\cite{masten2020inference}). To evaluate empirical 30-test FWER non-circularly under the least-favourable boundary null ($\Delta^L_{hl} = 0$), we conducted both a \textbf{stratified outbreak double bootstrap}~\cite{beran1988prepivoting} ($400$ outer worlds $\times 499$ inner resamples preserving $n_{js}$; no state correlation simulated) and a \textbf{53-cluster Bayesian double bootstrap} ($300$ outer $\text{Exp}(1)$ cluster-weighted worlds $\times 399$ inner draws simulating exposure-state and Multistate dependence), re-estimating every procedure's critical value inside each outer world. Finally, we evaluated 30-test inference under R2, R10, R4a/R4b, diagnostic-informed asymmetric bands, and exact leave-one-outbreak-out (LOO, all 3,176 outbreaks) and leave-largest-outbreak-out (LLO) deletions, alongside exploratory non-response follow-up simulations~\cite{hansen1946problem} (Supplementary Section S2, Tables S5--S11).
